% BEGIN REV09_RECTAS27_PROLONGACION
\subsubsection{Prolongación reversible en coordenadas de rectas}
\label{corpus9:xii:rectas-prolongacion}

La proposición~\ref{prop:cobertura-729-seccion} ha construido las
biyecciones \(\beta_{{\rm blk},N}\) mediante segmentos nonádicos del
registro ordenado, junto con su compatibilidad con la supresión de
bloques. El homeomorfismo
\eqref{eq:homeomorfismo-calibrado-ch} identifica el límite de esas
historias generadas con la sucesión de sus palabras de seis trits.
Componemos ahora ese lector con la coordenatización finita
\eqref{corpus9:xii:rectas-codificacion}, conservando su ley central.

\begin{theorem}[Transporte cilíndrico de historias a pares de rectas]
\label{corpus9:xii:rectas-homeomorfismo}
Para cada \(N\geq0\), la aplicación
\begin{equation}
 \mathcal E_N:=\mathcal E^{\times N}\circ\beta_{{\rm blk},N}:
 Y_N\xrightarrow{\ \cong\ }
 (\mathscr L_{27}\times\mathscr L_{27})^N
 \label{corpus9:xii:rectas-nivel}
\end{equation}
es una biyección. Estas biyecciones satisfacen
\begin{equation}
 \operatorname{pref}_N\circ\mathcal E_{N+1}
 =\mathcal E_N\circ\rho^{\rm blk}_{N+1,N}
 \label{corpus9:xii:rectas-truncamiento}
\end{equation}
y determinan el homeomorfismo
\begin{equation}
 \mathcal E_\infty:=\mathcal E^{\mathbb N}\circ\beta_{\rm blk}:
 \Omega_\Gamma^{\rm cal}\xrightarrow{\ \cong\ }
 (\mathscr L_{27}\times\mathscr L_{27})^{\mathbb N}.
 \label{corpus9:xii:rectas-limite}
\end{equation}
Su inversa es
\(\beta_{\rm blk}^{-1}\circ(\mathcal E^{-1})^{\mathbb N}\).
\end{theorem}
\begin{proof}
La proposición~\ref{corpus9:xii:rectas-biyeccion} proporciona la
inversa efectiva de cada factor \(\mathcal E\), mediante la tabla
de rectas y el entrelazado de coordenadas. Por ello
\(\mathcal E_N^{-1}=\beta_{{\rm blk},N}^{-1}
\circ(\mathcal E^{-1})^{\times N}\). Para \(N=0\), ambas aplicaciones
actúan sobre la palabra vacía. La aplicación coordenada a coordenada
conmuta con la supresión del último factor. Usando
\eqref{eq:cobertura-729-ch}, se obtiene
\[
 \begin{aligned}
 \operatorname{pref}_N\mathcal E^{\times(N+1)}
           \beta_{{\rm blk},N+1}
 &=\mathcal E^{\times N}\operatorname{pref}_N
           \beta_{{\rm blk},N+1}\\
 &=\mathcal E^{\times N}\beta_{{\rm blk},N}
           \rho^{\rm blk}_{N+1,N}.
 \end{aligned}
\]
Esta es \eqref{corpus9:xii:rectas-truncamiento}. Una sucesión de pares
de rectas determina, por la inversa de cada factor, una sucesión única
de palabras. El homeomorfismo
\eqref{eq:homeomorfismo-calibrado-ch} reconstruye a partir de ella una
historia única en \(\Omega_\Gamma^{\rm cal}\). Las composiciones
recuperan respectivamente la historia y la sucesión iniciales.
En las topologías cilíndricas, un abierto básico fija un número finito
de coordenadas. La aplicación y su inversa transforman esas condiciones
en condiciones sobre el mismo número de coordenadas; ambas son
continuas. Esto demuestra la afirmación topológica y la compatibilidad
del límite con todos los truncamientos finitos.
\end{proof}

Para conservar la información del estado completo, escribimos
\(\zeta\) para sus registros restantes: marco, orientación, acarreo,
ruta, memoria y datos de frontera. Sobre el dominio admisible ya
construido, la transformación ampliada es
\[
 \widehat{\mathcal E}(w,\zeta)=(\mathcal E(w),\zeta).
\]
Su dominio imagen conserva las mismas condiciones de compatibilidad,
expresadas mediante la inversa explícita de la primera coordenada. Si
\(F:\mathcal X\to\mathcal Y\) es un operador entre dominios de estados,
su transporte es
\[
 F^{\mathcal E}
 =\widehat{\mathcal E}_{\mathcal Y}\circ F
       \circ\widehat{\mathcal E}_{\mathcal X}^{-1}.
\]
La identidad \(F^{\mathcal E}\widehat{\mathcal E}_{\mathcal X}
=\widehat{\mathcal E}_{\mathcal Y}F\) se obtiene cancelando las dos
aplicaciones inversas. Para una prolongación relacional
\(\mathcal R\subset\mathcal X\times\mathcal Y\), la definición
correspondiente es
\[
 \mathcal R^{\mathcal E}
 =\{(\widehat{\mathcal E}_{\mathcal X}x,
       \widehat{\mathcal E}_{\mathcal Y}y):(x,y)\in\mathcal R\}.
\]
La transformación inversa recupera exactamente \(\mathcal R\).
Si la prolongación conserva además un testigo \(\omega\) en un dominio
\(W\), su realización \(\widetilde{\mathcal R}\subset
\mathcal X\times W\times\mathcal Y\) se transporta por
\[
 (x,\omega,y)\longmapsto
 (\widehat{\mathcal E}_{\mathcal X}x,\omega,
       \widehat{\mathcal E}_{\mathcal Y}y).
\]
La inversa actúa sobre los dos extremos y conserva \(\omega\). Así,
testigos diferentes siguen siendo diferentes, aunque compartan los
extremos, y la multiplicidad registrada permanece. La condición de
funcionalidad se conserva en los dominios en que ya estaba establecida.

La cobertura a profundidad arbitraria procede aquí del dominio
\(\Omega_\Gamma^{\rm cal}\) y de su lector probado, mientras los
pares de rectas proporcionan coordenadas reversibles de ese dominio.
El censo inicial de \(243\) palabras y la capacidad global de
\(729\) prolongaciones por bloque conservan sus funciones distintas.
La clasificación diagonal, incidente y alabeada acompaña cada factor;
las condiciones que seleccionan paseos adyacentes se transportan como
restricciones de la sucesión completa.
% END REV09_RECTAS27_PROLONGACION
